% Standalone problem document, generated from the adjacent README.md.
% Compile with XeLaTeX. Mathematical content is maintained in Markdown.
\documentclass[11pt,a4paper]{article}
\usepackage[fontsize=10.5pt]{scrextend}
\usepackage[margin=22mm,headheight=15pt,headsep=9mm,footskip=11mm]{geometry}
\usepackage{fontspec}
\setmainfont{texgyrepagella}[Extension=.otf,UprightFont=*-regular,BoldFont=*-bold,ItalicFont=*-italic,BoldItalicFont=*-bolditalic]
\setsansfont{texgyreheros}[Extension=.otf,UprightFont=*-regular,BoldFont=*-bold,ItalicFont=*-italic,BoldItalicFont=*-bolditalic]
\setmonofont{texgyrecursor}[Extension=.otf,UprightFont=*-regular,BoldFont=*-bold,ItalicFont=*-italic,BoldItalicFont=*-bolditalic,Scale=MatchLowercase]
\usepackage{amsmath,amssymb}
\usepackage{unicode-math}
\setmathfont{texgyrepagella-math.otf}
\usepackage{xcolor,microtype,enumitem,needspace,fancyhdr}
\usepackage{bookmark}
\definecolor{ink}{HTML}{17344B}
\definecolor{accent}{HTML}{197D80}
\definecolor{muted}{HTML}{596773}
\hypersetup{colorlinks=true,linkcolor=accent,urlcolor=accent,pdftitle={TR-28 - Entropy
formula for the asymptotic subrank of Dicke
tensors},pdfauthor={Open Problems in Numerical Linear Algebra}}
\urlstyle{same}
\setlength{\parindent}{0pt}
\setlength{\parskip}{5pt plus 1pt minus 1pt}
\linespread{1.04}
\setlength{\emergencystretch}{3em}
\widowpenalty=10000
\clubpenalty=10000
\displaywidowpenalty=10000
\allowdisplaybreaks[1]
\setlist{nosep,leftmargin=1.5em}
\providecommand{\tightlist}{\setlength{\itemsep}{0pt}\setlength{\parskip}{0pt}}
\providecommand{\pandocbounded}[1]{#1}
\pagestyle{fancy}
\fancyhf{}
\fancyhead[L]{\sffamily\footnotesize\color{muted}OPEN PROBLEMS IN NUMERICAL LINEAR ALGEBRA}
\fancyhead[R]{\sffamily\bfseries\footnotesize\color{accent}TR-28}
\fancyfoot[L]{\parbox[b]{0.9\textwidth}{\sffamily\fontsize{8}{10}\selectfont\color{muted}Editorial ratings; a bounded search cannot certify that a problem remains unsolved.\\Literature check: 2026-09-10}}
\fancyfoot[R]{\sffamily\footnotesize\color{muted}\thepage}
\renewcommand{\headrulewidth}{0.3pt}
\renewcommand{\headrule}{\hbox to\headwidth{\color{accent}\leaders\hrule height \headrulewidth\hfill}}
\makeatletter
\renewcommand{\section}{\Needspace{5\baselineskip}\@startsection{section}{1}{0pt}{12pt plus 2pt minus 2pt}{4pt}{\sffamily\bfseries\large\color{ink}}}
\renewcommand{\subsection}{\Needspace{4\baselineskip}\@startsection{subsection}{2}{0pt}{9pt plus 1pt minus 1pt}{3pt}{\sffamily\bfseries\normalsize\color{ink}}}
\makeatother
\setcounter{secnumdepth}{0}
\begin{document}
{\sffamily\small\color{muted}Tensor computations\par}
\vspace{3pt}
{\raggedright\hyphenpenalty=10000\exhyphenpenalty=10000\sffamily\bfseries\fontsize{21}{24}\selectfont\color{ink}Entropy
formula for the asymptotic subrank of Dicke tensors\par}
\vspace{7pt}
{\sffamily\small\textbf{Difficulty:} challenging\quad\textbf{Importance:} interesting
to the community\par}
{\sffamily\small\bfseries\color{accent}Status: Partially resolved\par}
\vspace{4pt}
{\color{accent}\hrule height 0.8pt}
\vspace{6pt}
\textbf{Rating rationale:} Matching the entropy upper bound for this
structured tensor family requires substantial progress beyond its known
binary and multilinear cases. Tensor restriction rates and entanglement
conversion give the question community importance.

\subsection{Problem statement}\label{problem-statement}

Let \(n\geq2\) and let \(d_1,\ldots,d_n\) be positive integers with
\(k=d_1+\cdots+d_n\). In \((\mathbb C^n)^{\otimes k}\), define the Dicke
tensor \[
D_{\mathbf d}=
\sum_{\substack{(i_1,\ldots,i_k)\in\{1,\ldots,n\}^{k}\\
|\{j:i_j=a\}|=d_a\ \text{for all }a}}
e_{i_1}\otimes\cdots\otimes e_{i_k}.
\] For any tensor \(T\in V_1\otimes\cdots\otimes V_k\), its subrank
\(Q(T)\) is the largest integer \(s\) for which there are linear maps
\(L_j:V_j\to\mathbb C^s\) satisfying \[
(L_1\otimes\cdots\otimes L_k)T=\sum_{a=1}^{s}e_a^{\otimes k}.
\] Let \(\widetilde Q(T)=\lim_{m\to\infty}Q(T^{\boxtimes m})^{1/m}\),
where powers group corresponding modes.

Does every Dicke tensor satisfy \[
\widetilde Q(D_{\mathbf d})
=2^{H(d_1/k,\ldots,d_n/k)},\qquad
H(p_1,\ldots,p_n)=-\sum_{a=1}^np_a\log_2p_a?
\] The maps in the subrank definition are arbitrary complex linear maps,
independently chosen in each mode. No symmetry or coordinate-selection
constraint is imposed. Multiplying \(D_{\mathbf d}\) by a nonzero scalar
does not change the question, so this agrees with the symmetric-tensor
convention for the monomial \(x_1^{d_1}\cdots x_n^{d_n}\).

\subsection{Why it matters}\label{why-it-matters}

Subrank measures how many independent scalar computations a tensor can
perform after linear transformations. The conjecture predicts the
exponential rate obtainable from repeated copies of an explicit
structured tensor solely from its multiplicity distribution.

\subsection{References and status
check}\label{references-and-status-check}

\begin{enumerate}
\def\labelenumi{\arabic{enumi}.}
\tightlist
\item
  P. Vrana and M. Christandl, \emph{Asymptotic entanglement
  transformation between W and GHZ states},
  \href{https://arxiv.org/pdf/1310.3244}{arXiv:1310.3244}, §VI, equation
  (40), p.10; Journal of Mathematical Physics \textbf{56} (2015),
  022204. The conversion-rate formulation is reciprocal to the base-two
  logarithm of asymptotic subrank.
\item
  A. Oneto and E. Ventura, \emph{Ranks of tensors: geometry and
  applications}, \href{https://doi.org/10.1007/s40574-025-00472-9}{2025
  survey}, Conjecture 4.26, states the displayed subrank formula
  explicitly.
\item
  M. Christandl, P. Vrana, and J. Zuiddam, \emph{Asymptotic tensor rank
  of graph tensors: beyond matrix multiplication},
  \href{https://ir.cwi.nl/pub/28609/28609.pdf}{Computational Complexity
  \textbf{28} (2019), 57--111}, introduction and tight-tensor
  lower-bound results.
\end{enumerate}

\subsubsection{Status check ---
2026-09-10}\label{status-check-2026-09-10}

Checked the source formulation and the explicit 2025 restatement,
Conjecture 4.26, with targeted Dicke-tensor/subrank/entropy searches
through 2026. The formula is established for \(n=2\) and for
\(d_1=\cdots=d_n=1\), substantive parameter families within this target.
The all-multiplicity statement remains conjectural in the survey, and no
full later resolution was located. Asymptotic-rank computations and
lower bounds on subrank do not by themselves establish the requested
equality.
\end{document}
